%% Chapter 16 ----------------------------------------------------------------
\chapter{The Solar Image Analysis Window}
\label{ch:solar-window}
\index{Solar Image Analysis}

Solar Image Analysis is a multi-mission imaging workspace for extreme-ultraviolet (EUV)
disk images, magnetograms, coronagraph images and heliospheric images. It retrieves image
sequences from the solar archives or reads local FITS files, displays and processes them,
measures eruptions and their kinematics, and models the coronal magnetic field\index{magnetic field}. Open it with
\menu{Analysis > Solar Image Analysis} in the main window.

This chapter describes the window, its menus and sessions, and how image data are obtained.
Display and processing are described in Chapter~\ref{ch:solar-processing}, measurements in
Chapter~\ref{ch:solar-measurements} and magnetic-field modeling in Chapter~\ref{ch:pfss}.

\section{Supported observables}
\label{sec:observables}
\index{observable}

Table~\ref{tab:observables} lists the observables that can be retrieved. Local FITS files
from these and other solar instruments can be loaded as well.

\begin{table}[!htbp]
  \centering
  \caption{Observables available in Solar Image Analysis.}
  \label{tab:observables}
  \small
  \begin{tabularx}{\linewidth}{@{}P{0.25\linewidth}P{0.18\linewidth}L@{}}
    \toprule
    \thead{Instrument} & \thead{Class} & \thead{Observables} \\
    \midrule
    SDO/AIA\index{SDO/AIA} \parencite{lemen2012} & EUV/UV disk & 94, 131, 171, 193, 211, 304, 335, 1600 and 1700~Å \\
    SDO/HMI\index{SDO/HMI} \parencite{scherrer2012} & Magnetograph & Magnetogram, Continuum (Intensitygram), Dopplergram; vector field (Section~\ref{sec:hmi-vector}) \\
    SOHO/EIT\index{SOHO/EIT} & EUV disk & 171, 195, 284 and 304~Å \\
    SOHO/LASCO\index{SOHO/LASCO} \parencite{brueckner1995} & Coronagraph & C2, C3 \\
    STEREO-A\index{STEREO}/B SECCHI \parencite{howard2008} & EUV disk, coronagraph, heliospheric imager & EUVI 171, 195, 284 and 304~Å; COR1, COR2; HI1, HI2 \\
    GOES/SUVI\index{GOES/SUVI} & EUV disk & 94, 131, 171, 195, 284 and 304~Å \\
    \bottomrule
  \end{tabularx}
\end{table}

The sidebar shows only the tool groups that apply to the instrument class of the loaded data
(or, before loading, of the selected observable); see Table~\ref{tab:solar-cards}.

\section{Window layout}
\label{sec:solar-layout}

The window is divided into the regions numbered in Figure~\ref{fig:solar-window} and
described in Table~\ref{tab:solar-regions}.

\screenshot{solar/solar_window_overview}{The Solar Image Analysis window with its regions numbered.}{fig:solar-window}{Solar Image Analysis window with an AIA 193 Å sequence loaded and Measurements enabled. Add numbered callouts: 1 menu bar, 2 sidebar cards, 3 sidebar handle, 4 coordinate readout and Export MP4, 5 Measure toolbar, 6 image canvas, 7 CME tracking panel, 8 playback bar, 9 Details drawer.}

\begin{table}[!htbp]
  \centering
  \caption{Regions of the Solar Image Analysis window.}
  \label{tab:solar-regions}
  \small
  \begin{tabularx}{\linewidth}{@{}cP{0.22\linewidth}L@{}}
    \toprule
    \thead{No.} & \thead{Region} & \thead{Description} \\
    \midrule
    \callout{1} & Menu bar & \menu{Session}, \menu{Data}, \menu{Analysis}, \menu{Movie}, \menu{Export} and \menu{About} (Section~\ref{sec:solar-menus}). \\
    \callout{2} & Sidebar & Collapsible cards that read from top to bottom as the workflow: get data, choose records, analyze, then display, export and instrument-specific tools. \\
    \callout{3} & Sidebar handle & Hides or shows the sidebar. \\
    \callout{4} & Header bar & Live cursor readout --- distance from Sun center in \(R_\odot\), position angle, heliographic coordinates on the disk, and pixel --- and the one-click \gui{Export MP4}\index{movie export!MP4} button. \\
    \callout{5} & Measure toolbar & The \gui{Measurements}\index{measurement tools!enabling} switch, the \gui{Ruler}\index{measurement tools!ruler}, \gui{Profile}\index{measurement tools!profile}, \gui{Region Stats}\index{measurement tools!region statistics}, \gui{Track CME}\index{Track CME}, \gui{Circle Fit}\index{circle fit} and \gui{Clear} tools, and \gui{Pan / Zoom}\index{Solar Image Analysis!pan and zoom} (Chapter~\ref{ch:solar-measurements}). \\
    \callout{6} & Image canvas & The current frame, titled with instrument, wavelength and time. \\
    \callout{7} & CME tracking panel & Measurement table, live height--time graph and kinematic fit; active when \gui{Measurements} is ticked (Section~\ref{sec:tracking-panel}). \\
    \callout{8} & Playback bar & Rewind, previous, play, pause and next buttons, the frame scrubber, the frame counter (\gui{Frame n / N}) and the \gui{Speed} in frames per second. \\
    \callout{9} & Details drawer\index{Solar Image Analysis!Details drawer} & Results and warnings of every operation. Click \gui{Details} to collapse it; drag the divider above it to resize. \\
    \bottomrule
  \end{tabularx}
\end{table}

The playback speed set in the playback bar is used for playback and for exported movies.
While a sequence plays, or while a contrast slider is dragged, the canvas switches to a fast
preview that subsamples the display; full resolution returns as soon as playback stops.
Measurements, statistics and exports always use the full-resolution data.

\subsection{Menus}
\label{sec:solar-menus}

\begin{xltabular}{\linewidth}{@{}P{0.38\linewidth} L@{}}
  \caption{Menus of the Solar Image Analysis window.}\label{tab:solar-menus}\\
  \toprule \thead{Command} & \thead{Function} \\ \midrule\endfirsthead
  \toprule \thead{Command} & \thead{Function} \\ \midrule\endhead
  \menu{Session > Open Session…} & Open an \file{.ecsolar}\index{ecsolar@.ecsolar file} session (Section~\ref{sec:solar-sessions}). \\
  \menu{Session > Save Session} / \menu{Save Session As…} & Save the session (\kbd{Ctrl+S} / \kbd{Ctrl+Shift+S}). \\
  \menu{Data > Fetch Archive Records} & Search the archive (\gui{Fetch Records}\index{Solar Image Analysis!Fetch Records}). \\
  \menu{Data > Find Latest Available}\index{Solar Image Analysis!Find Latest} & Find the newest available data (\gui{Find Latest}). \\
  \menu{Data > SOHO Live Preview (Helioviewer)}\index{Helioviewer} & Near-real-time SOHO quicklook images (Section~\ref{sec:live-preview}). \\
  \menu{Data > Load Selected (to cache)}\index{Solar Image Analysis!Load Selected} & Download and load the ticked records. \\
  \menu{Data > Save Selected to Disk…} & Download the ticked records to a folder. \\
  \menu{Data > Upload FITS Files}\index{Solar Image Analysis!local files} & Load local FITS files. \\
  \menu{Data > Use Analyzer Window} & Copy the time range of the main window's spectrum. \\
  \menu{Data > Stop Download/Search} & Cancel the running operation. \\
  \menu{Data > Reset All (clear cache \& defaults)} & Clear the cache and restore default settings. \\
  \menu{Analysis > Plot Frames}, \menu{Running Difference}\index{running difference}, \menu{Composite} & Display modes (Section~\ref{sec:solar-analysis-card}). \\
  \menu{Analysis > Identify Active Regions}, \menu{Fetch NOAA/HEK Labels}\index{active regions!NOAA labels} & Active-region tools (Section~\ref{sec:active-regions}). \\
  \menu{Analysis > Compare Viewpoint…} & Compare with a second observer (Section~\ref{sec:compare-viewpoint}). \\
  \menu{Analysis > GCS CME Fitting…} & Open GCS fitting at the time of the displayed frame (Chapter~\ref{ch:gcs-window}). \\
  \menu{Analysis > Reset Loaded Frames} & Restore the frames as loaded. \\
  \menu{Movie > Rewind}, \menu{Back}, \menu{Play}, \menu{Pause}, \menu{Forward} & Playback. \\
  \menu{Movie > Build Movie}\index{movie export} & Export a movie (Section~\ref{sec:solar-movies}). \\
  \menu{Export > Export Plot}, \menu{Export Cropped FITS}\index{crop!export FITS}, \menu{Export Regions CSV}, \menu{Export MP4} & Exports (Section~\ref{sec:solar-movies}). \\
  \menu{About > About Solar Image Analysis…} & Version information. \\
  \bottomrule
\end{xltabular}

\subsection{Sidebar cards}

Table~\ref{tab:solar-cards} lists the sidebar cards in order, and the instrument classes for
which each is shown.

\begin{table}[!htbp]
  \centering
  \caption{Sidebar cards of Solar Image Analysis.}
  \label{tab:solar-cards}
  \small
  \begin{tabularx}{\linewidth}{@{}P{0.3\linewidth}P{0.3\linewidth}L@{}}
    \toprule
    \thead{Card} & \thead{Shown for} & \thead{See} \\
    \midrule
    Data Source & all & Section~\ref{sec:data-source} \\
    Archive Results & all & Section~\ref{sec:archive-results} \\
    Analysis & all (disk tools for disk imagers and magnetographs only) & Section~\ref{sec:solar-analysis-card} \\
    Processing Level & AIA and SUVI & Section~\ref{sec:processing-level} \\
    Display \& Crop & all & Section~\ref{sec:display-crop} \\
    Solar Derotation & all & Section~\ref{sec:derotation} \\
    Movie Export & all & Section~\ref{sec:solar-movies} \\
    Coronagraph Tools & coronagraphs & Section~\ref{sec:nrgf} \\
    Overlay Layers & coronagraphs & Section~\ref{sec:overlay-layers} \\
    Heliospheric Imager (J-map) & HI1, HI2 & Section~\ref{sec:jmap} \\
    Magnetic Vector Field (HMI) & HMI & Section~\ref{sec:hmi-vector} \\
    Potential Field (PFSS) & disk imagers, magnetographs, synoptic maps & Chapter~\ref{ch:pfss} \\
    Active Regions & disk imagers, magnetographs & Section~\ref{sec:active-regions} \\
    \bottomrule
  \end{tabularx}
\end{table}

\section{Obtaining image data}

\subsection{The Data Source card}
\label{sec:data-source}

The \gui{Data Source} card (Figure~\ref{fig:solar-data-source}) defines what to retrieve and
how.

\screenshot[0.5\linewidth]{solar/solar_data_source}{The \gui{Data Source} card.}{fig:solar-data-source}{The Data Source card (Observable, Start/End, Cadence, Max records, Use Analyzer Time Window, Download options, buttons) above the Archive Results card with a list of records, some ticked.}

\begin{description}
  \item[Observable] the instrument and wavelength or product (Table~\ref{tab:observables}).
  \item[Start (UTC), End (UTC)] the time range.
  \item[Cadence (s)] the sampling interval between frames (0--3600~s; default 120~s).
  \item[Max records] the maximum number of records to list (1--5000; default 120).
  \item[Use Analyzer Time Window] copies the start and end times of the spectrum in the main
    window, so that the images match the radio burst being studied.
\end{description}

\procedure{Download options (SDO)}
\begin{description}
  \item[Source] \gui{Auto (JSOC → VSO)}\index{VSO}\index{JSOC} tries the Joint Science Operations Center (JSOC)
    first and falls back to the Virtual Solar Observatory (VSO); \gui{JSOC (fast)} returns
    direct, compressed AIA files and is usually much faster; \gui{VSO (classic)} uses the
    VSO only.
  \item[JSOC notify e-mail] JSOC requires a one-time, free registration of an e-mail address
    at \url{https://jsoc.stanford.edu/ajax/register_email.html}. The address is remembered.
  \item[Frame size] \gui{Full disk (4096²)}, \gui{Binned ½ (2048²)}, \gui{Binned ¼ (1024²)}
    or \gui{Cutout (region)}\index{Solar Image Analysis!cutout}. Binned and cutout frames are produced by JSOC, so only the
    reduced data are transferred; they require the JSOC source and e-mail. For a cutout,
    enter the \gui{Centre X″}, \gui{Centre Y″}, \gui{Width″} and \gui{Height″} in arcseconds
    from disk center.
  \item[High resolution AIA (best crop)] requests full-resolution AIA products. Files are
    larger, but cropped views keep more detail.
\end{description}
An estimate of the download size is shown below the options when records are selected.

\procedure{Actions}
\begin{description}
  \item[Fetch Records] searches the archive; the matching records appear in
    \gui{Archive Results}. \gui{Stop} cancels a running search or download; files already
    completed remain in the cache.
  \item[Find Latest] scans the archive backwards to find the newest data actually available
    for the observable. This is useful for SOHO/LASCO, whose calibrated archive can lag real
    time by many months, and may take a minute or two.
  \item[Upload FITS…] loads solar FITS files from disk without an archive search.
  \item[Live Preview (Helioviewer)] shown for SOHO/LASCO and SOHO/EIT
    (Section~\ref{sec:live-preview}).
\end{description}

\subsection{The Archive Results card}
\label{sec:archive-results}

The \gui{Archive Results} table lists the records found, with the columns \gui{Use},
\gui{Start UTC}, \gui{Source}, \gui{Size} and \gui{File ID}. Tick records in the \gui{Use}
column, or use \gui{Select All} and \gui{Deselect All}, then:
\begin{itemize}
  \item \gui{Load Selected} downloads the ticked records into the working cache and loads
    them for analysis; or
  \item \gui{Save to Disk…} downloads them to a folder you choose, for keeping.
\end{itemize}
The \gui{Analysis} card then summarizes what is loaded (instrument, number of frames and
configuration).

\begin{note}[SOHO data]
  The archive serves two files for every SOHO/EIT observation, the calibrated Level~1 file
  and the raw level-zero file. The window chooses one per observation: Level~1 where it
  exists and the raw file otherwise. Level~1 processing runs about a year behind the raw
  archive, so recent searches return raw frames and older ones calibrated frames; a window
  that straddles the boundary loads as two configuration groups. The calibrated SOHO/LASCO
  archive also lags real time. For a date beyond the newest archived data, the search falls
  back to the nearest available data and says so. Use \gui{Live Preview} to see the most
  recent SOHO images.
\end{note}

\subsection{Live Preview (Helioviewer)}
\label{sec:live-preview}
\index{Helioviewer}

\gui{Live Preview (Helioviewer)} opens the \gui{SOHO/… Near-Real-Time Preview (Helioviewer)}
window (Figure~\ref{fig:live-preview}) with the newest quicklook image of SOHO/LASCO C2 or C3
or SOHO/EIT, usually updated within about an hour. These browse images are for situational
awareness, not for analysis.

\screenshot[0.5\linewidth]{solar/helioviewer_preview}{The SOHO near-real-time preview.}{fig:live-preview}{The SOHO/LASCO Near-Real-Time Preview (Helioviewer) window showing the latest C2 image, the Detector selector, the Latest Frame button, and the Movie (time range · step) group with playback controls.}

\begin{description}
  \item[Passband / Detector, Latest Frame] choose the EIT passband or LASCO detector, and
    reload the newest image.
  \item[Movie (time range · step)] set \gui{Start (UTC)}, \gui{End (UTC)}, \gui{Step}
    (minutes) and \gui{Max frames}, then \gui{Build Movie} to fetch one frame per step and
    play them with \gui{◀}, \gui{Play}, \gui{Pause}, \gui{▶} and an FPS setting.
    \gui{Cancel} stops the fetch. The frame cap protects the Helioviewer service.
  \item[Open in Browser] opens the view on the Helioviewer web site.
\end{description}

\subsection{Processing level}
\label{sec:processing-level}
\index{calibration level}

Downloads always fetch the archive's base product. For instruments that publish more than
one processing level, the \gui{Processing Level} card selects the level that is displayed and
analyzed:
\begin{description}
  \item[SDO/AIA] \gui{Level 1 (as downloaded)}, the archive product; or
    \gui{Level 1.5 (registered here)}\index{calibration level!AIA level 1.5}, computed locally with aiapy \parencite{barnes2020},
    which applies the master pointing, rotates the image to solar north, rescales it to
    0.6~arcsec per pixel and centers the Sun.
  \item[GOES/SUVI] \gui{Level 1b (as downloaded)}, single calibrated exposures; or
    \gui{Level 2 (archive composite)}, NOAA's high-dynamic-range composites, which are
    downloaded separately when selected.
\end{description}
\gui{Apply Level}\index{calibration level!Apply Level} re-derives the loaded sequence at the selected level; any crop or
derotation is re-applied afterwards. The status line reports the level in use. The card is
hidden for instruments that offer only one level.

\section{Sessions}
\label{sec:solar-sessions}
\index{session!solar}\index{ecsolar@.ecsolar file}

A Solar Image Analysis session (\file{.ecsolar}) is a single file that reopens an analysis
exactly where it was left. It embeds the original FITS frames, so it survives a cleared
cache or a move to another computer (at the cost of a large file), together with the data
source, the display, crop and playback state, the CME measurements (height--time picks,
circle fits and fit order) and the PFSS settings. The PFSS solution itself is not stored; one
click on \gui{Compute PFSS}\index{PFSS!compute} reproduces it. Sessions are opened and saved from the
\menu{Session} menu (\kbd{Ctrl+S}, \kbd{Ctrl+Shift+S}).
